\textbf{Цели работы:}
\begin{itemize}
	\item изучить базовое устройство Linux;
	\item научиться взаимодействовать с ОС Linux;
	\item научиться писать bash-скрипты.
\end{itemize}

\textbf{Задачи:}
\begin{enumerate}
	\item изучить методические материалы к лабораторной работе; 
	\item написать bash-скрипт для удаления файлов по заданному имени из указанной директории, а также предусмотреть флаг -p для удаления по шаблону. 
\end{enumerate}

\section*{Выполнение задания}

\begin{figure}[!htbp]
	\centering
	\includegraphics[width=0.6\linewidth]{"Screenshot (53)"}
	\caption{Установленная ОС Ubuntu}
\end{figure}
\FloatBarrier

\begin{figure}[!htbp]
	\centering
	\includegraphics[width=0.6\linewidth]{"Screenshot (54)"}
	\caption{Терминал}
\end{figure}
\FloatBarrier

\clearpage

\section*{Исходный код программы}
% согласно стилю языка программирования
% consolas, 10

\lstinputlisting[style=bash]{deleteFiles.sh}
%\lstinputlisting[style=bash, linerange={}]{}
Листинг 1 -- bash-скрипт по заданию

\begin{figure}[!htbp]
	\centering
	\includegraphics[width=0.6\linewidth]{"Screenshot From 2026-02-26 15-33-52"}
	\caption{Вывод программы}
\end{figure}
\FloatBarrier

\section*{Заключение}

В ходе выполнения лабораторной работы были решены следующие задачи:
\begin{itemize}
	\item изучено базовое устройство Linux;
	\item получены знания о том, как взаимодействовать с ОС Linux;
	\item получены знания о том, как писать bash-скрипты.
\end{itemize}

\section*{Контрольные вопросы}

\begin{enumerate}
	\item \textit{Что такое Linux?} Компьютерная операционная система, разработанная на основе модели открытого исходного кода.
	\item \textit{В чем отличие переменных окружения от переменных оболочки?} Операционная система передает все переменные окружения системным процессам и программам, выполняемым оболочкой, в то время как переменные оболочки доступны только внутри текущего экземпляра оболочки.
	\item \textit{С помощью какой команды изменяются права доступа к файлу?} Командой chmod. % chmod u+rwx,o-wx file
\end{enumerate}
